\section{The \texttt{Grating\_trans} McXtrace Component}
Transmission grating

\subsection*{Identification}
\begin{itemize}
  \item \textbf{Author:} Erik B Knudsen (erkn@fysik.dtu.dk)
  \item \textbf{Origin:} DTU Physics
  \item \textbf{Date:} December 2016
\end{itemize}

\subsection*{Description}
Model of a 1D rectangular transmission grating based on the theory developed in Schnopper et. al., Applied Optics, 1977. The grating lines are assumed to be vertical. Within each period a fraction gamma is the "open" fraction. (I.e. 1 is completely open). At present only absorption in the substrate (modelled by the thickness sdepth) is included.

This component is currently undergoing validation.

Example: Grating\_trans( xwidth=25e-3, yheight=25e-3, gamma=0.4, period=2000e-10, zdepth=5100e-10, max\_order=3, material="Au.txt")

\subsection*{Input parameters}
Parameters in \textbf{boldface} are required; the others are optional.

\begin{longtable}{p{0.22\textwidth}p{0.12\textwidth}p{0.46\textwidth}p{0.14\textwidth}}
\toprule
\textbf{Name} & \textbf{Unit} & \textbf{Description} & \textbf{Default} \\
\midrule
\endhead
xwidth & m & Width of the grating. Defines how many lines there are in total. & 1e-3 \\
yheight & m & Height of the grating. & 1e-3 \\
period & m & Distance between grating grooves. & 1e-6 \\
gamma & 0-1 & Ratio between groove and period  aka duty cycle. 1 means fully open. & 0.5 \\
zdepth & m & Depth of grooves. & 1e-6 \\
sdepth & m & Thickness of substrate. The default is to have no substrate - i.e. rods. & 0 \\
material & str & Data file containing the material from which the grating is made. & "Au.txt" \\
substrate & str & Data file containing material data for the substrate. & "" \\
max\_order & 1 & Maximum order to diffract & 2 \\
fixed\_delta & 0/1 & Set delta to the given constant. Useful for debugging. & 0 \\
\bottomrule
\end{longtable}

\subsection*{Links}
\begin{itemize}
  \item Component source code found in file \texttt{Grating\_trans.comp}.
\end{itemize}
\IfFileExists{optics/Grating_trans_static.tex}{\input{optics/Grating_trans_static.tex}}{}